\documentclass[10pt,a4paper]{amsart}
\usepackage{amsmath,amssymb,amsthm,graphicx}
\usepackage{verbatim}
\newenvironment{reference}{\comment}{\endcomment}
%\newenvironment{reference}{}{}
\newenvironment{slogan}{\comment}{\endcomment}
\newenvironment{history}{\comment}{\endcomment}

% For commutative diagrams we use Xy-pic
\usepackage[all]{xy}

% We use 2cell for 2-commutative diagrams.
\xyoption{2cell}
% Theorem environments.
%
\theoremstyle{plain}
\newtheorem{theorem}[subsection]{Theorem}
\newtheorem{proposition}[subsection]{Proposition}
\newtheorem{lemma}[subsection]{Lemma}

\theoremstyle{definition}
\newtheorem{definition}[subsection]{Definition}
\newtheorem{example}[subsection]{Example}
\newtheorem{exercise}[subsection]{Exercise}
\newtheorem{situation}[subsection]{Situation}

\theoremstyle{remark}
\newtheorem{remark}[subsection]{Remark}
\newtheorem{remarks}[subsection]{Remarks}

\numberwithin{equation}{subsection}

% Macros
%
\def\lim{\mathop{\mathrm{lim}}\nolimits}
\def\colim{\mathop{\mathrm{colim}}\nolimits}
\def\Spec{\mathop{\mathrm{Spec}}}
\def\Hom{\mathop{\mathrm{Hom}}\nolimits}
\def\Ext{\mathop{\mathrm{Ext}}\nolimits}
\def\SheafHom{\mathop{\mathcal{H}\!\mathit{om}}\nolimits}
\def\SheafExt{\mathop{\mathcal{E}\!\mathit{xt}}\nolimits}
\def\Sch{\mathit{Sch}}
\def\Mor{\mathop{\mathrm{Mor}}\nolimits}
\def\Ob{\mathop{\mathrm{Ob}}\nolimits}
\def\Sh{\mathop{\mathit{Sh}}\nolimits}
\def\NL{\mathop{N\!L}\nolimits}
\def\CH{\mathop{\mathrm{CH}}\nolimits}
\def\proetale{{pro\text{-}\acute{e}tale}}
\def\etale{{\acute{e}tale}}
\def\QCoh{\mathit{QCoh}}
\def\Ker{\mathop{\mathrm{Ker}}}
\def\Im{\mathop{\mathrm{Im}}}
\def\Coker{\mathop{\mathrm{Coker}}}
\def\Coim{\mathop{\mathrm{Coim}}}

% Boxtimes
%
\DeclareMathSymbol{\boxtimes}{\mathbin}{AMSa}{"02}

%
% Macros for moduli stacks/spaces
%
\def\QCohstack{\mathcal{QC}\!\mathit{oh}}
\def\Cohstack{\mathcal{C}\!\mathit{oh}}
\def\Spacesstack{\mathcal{S}\!\mathit{paces}}
\def\Quotfunctor{\mathrm{Quot}}
\def\Hilbfunctor{\mathrm{Hilb}}
\def\Curvesstack{\mathcal{C}\!\mathit{urves}}
\def\Polarizedstack{\mathcal{P}\!\mathit{olarized}}
\def\Complexesstack{\mathcal{C}\!\mathit{omplexes}}
% \Pic is the operator that assigns to X its picard group, usage \Pic(X)
% \Picardstack_{X/B} denotes the Picard stack of X over B
% \Picardfunctor_{X/B} denotes the Picard functor of X over B
\def\Pic{\mathop{\mathrm{Pic}}\nolimits}
\def\Picardstack{\mathcal{P}\!\mathit{ic}}
\def\Picardfunctor{\mathrm{Pic}}
\def\Deformationcategory{\mathcal{D}\!\mathit{ef}}
\usepackage{cleveref}
\crefname{originallemma}{original lemma}{original lemmas}
\Crefname{originallemma}{Original lemma}{Original lemmas}
\makeatletter
\newlabel{perfect-lemma-affine-compare-bounded}{{06Z0}{source}{}{sourcepointer-1}{}}
\newlabel{perfect-lemma-affine-compare-bounded@cref}{{[originallemma][0][]06Z0}{[1][0][]source}{}{}{}}
\makeatother
\crefname{originallemma}{original lemma}{original lemmas}
\Crefname{originallemma}{Original lemma}{Original lemmas}
\makeatletter
\newlabel{schemes-lemma-standard-open-two-affines}{{01IW}{source}{}{sourcepointer-2}{}}
\newlabel{schemes-lemma-standard-open-two-affines@cref}{{[originallemma][0][]01IW}{[1][0][]source}{}{}{}}
\makeatother
\crefname{originallemma}{original lemma}{original lemmas}
\Crefname{originallemma}{Original lemma}{Original lemmas}
\makeatletter
\newlabel{dualizing-lemma-dualizing-localize}{{0A7G}{source}{}{sourcepointer-3}{}}
\newlabel{dualizing-lemma-dualizing-localize@cref}{{[originallemma][0][]0A7G}{[1][0][]source}{}{}{}}
\makeatother
\crefname{originallemma}{original lemma}{original lemmas}
\Crefname{originallemma}{Original lemma}{Original lemmas}
\makeatletter
\newlabel{dualizing-lemma-dualizing-glue}{{0A7H}{source}{}{sourcepointer-4}{}}
\newlabel{dualizing-lemma-dualizing-glue@cref}{{[originallemma][0][]0A7H}{[1][0][]source}{}{}{}}
\makeatother
\title{Being locally the image of a dualizing complex is independent of the affine open covering}
\author{The Stacks Project Authors}
\date{stacks-project:duality.tex:local-snapshot-20261012; GNU FDL 1.2 or later}
\begin{document}
\maketitle
\noindent\textit{Editorial note.} All delivered mathematical text is the author's own, copied from the pinned Stacks Project chapter file duality.tex; no statement or proof step was rewritten, reordered or removed.\par\medskip
\noindent\textit{Editorial note.} Preamble provenance: the pinned chapter file opens with an input of the shared preamble and carries no preamble text of its own, so the project's single Stacks-wide preamble staged at sources/stacks-duality/preamble.tex is excerpted verbatim under the recipe fixed by LANE\_BRIEF 22.1, as three spans, lines 8-13 (verbatim package and the reference/slogan/history environments), lines 14-18 (Xy-pic setup) and lines 154-217 (the native theorem declarations and mathematical macros). Nothing was added, removed or redefined; the wrapper class and the amsmath/amssymb/amsthm/graphicx base packages are supplied by the assembler.\par\medskip
\noindent\textit{Editorial note.} The counted statement and proof contain no cross reference, no citation and no included file, so the delivered body is self contained.\par\medskip
\noindent\textit{Editorial note.} Canonical location of the counted proposition: Stacks tag 0A86 (sources/stacks-duality/pinned-original-source.tex line 308).\par\medskip
\noindent\textit{Allowed standard prerequisites (editorial).}\begin{itemize}
\item Comparison of complexes on an affine scheme with complexes of modules. Derived Categories of Schemes, Lemma labelled perfect-lemma-affine-compare-bounded, cited at proof lines 324-328 and 340-344. Produces omega\_A and identifies the localisations.
\item Standard open coverings subordinate to a given affine cover. Schemes, Lemma labelled schemes-lemma-standard-open-two-affines, cited at proof lines 333-337. Provides the D(f\_j) inside the U\_i.
\item Localisation and glueing of dualizing complexes. Dualizing Complexes, Lemmas labelled dualizing-lemma-dualizing-localize and dualizing-lemma-dualizing-glue, cited at proof lines 348-352. Give dualizingness on the localisations and then for omega\_A.
\end{itemize}\medskip
\ifdefined\hypertarget\hypertarget{sourcepointer-1}{}\fi
\noindent\textit{Source reference.} Original lemma 06Z0: sources/stacks-chow/tags.txt, source lines 8904--8904. Cross-chapter standard prerequisite located in the project tag list staged in the local snapshot; invoked at proof line 329 of the counted proof. It is retained only as a navigation pointer; every load-bearing step is inside the delivered span.\par\medskip
\ifdefined\hypertarget\hypertarget{sourcepointer-2}{}\fi
\noindent\textit{Source reference.} Original lemma 01IW: sources/stacks-chow/tags.txt, source lines 1993--1993. Cross-chapter standard prerequisite located in the project tag list staged in the local snapshot; invoked at proof line 337 of the counted proof. It is retained only as a navigation pointer; every load-bearing step is inside the delivered span.\par\medskip
\ifdefined\hypertarget\hypertarget{sourcepointer-3}{}\fi
\noindent\textit{Source reference.} Original lemma 0A7G: sources/stacks-chow/tags.txt, source lines 12879--12879. Cross-chapter standard prerequisite located in the project tag list staged in the local snapshot; invoked at proof line 346 of the counted proof. It is retained only as a navigation pointer; every load-bearing step is inside the delivered span.\par\medskip
\ifdefined\hypertarget\hypertarget{sourcepointer-4}{}\fi
\noindent\textit{Source reference.} Original lemma 0A7H: sources/stacks-chow/tags.txt, source lines 12880--12880. Cross-chapter standard prerequisite located in the project tag list staged in the local snapshot; invoked at proof line 351 of the counted proof. It is retained only as a navigation pointer; every load-bearing step is inside the delivered span.\par\medskip
\begin{lemma}
\label{lemma-equivalent-definitions}
Let $X$ be a locally Noetherian scheme. Let $K$ be an object of
$D(\mathcal{O}_X)$. The following are equivalent
\begin{enumerate}
\item For every affine open $U = \Spec(A) \subset X$ there exists
a dualizing complex $\omega_A^\bullet$ for $A$ such that
$K|_U$ is isomorphic to the image of $\omega_A^\bullet$ by
the functor $\widetilde{} : D(A) \to D(\mathcal{O}_U)$.
\item There is an affine open covering $X = \bigcup U_i$, $U_i = \Spec(A_i)$
such that for each $i$ there exists a dualizing complex $\omega_i^\bullet$ for
$A_i$ such that $K|_{U_i}$ is isomorphic to the image of $\omega_i^\bullet$ by
the functor $\widetilde{} : D(A_i) \to D(\mathcal{O}_{U_i})$.
\end{enumerate}
\end{lemma}

\begin{proof}
Assume (2) and let $U = \Spec(A)$ be an affine open of $X$.
Since condition (2) implies that $K$ is in $D_\QCoh(\mathcal{O}_X)$
we find an object $\omega_A^\bullet$ in $D(A)$ whose associated
complex of quasi-coherent sheaves is isomorphic to $K|_U$, see
Derived Categories of Schemes, Lemma
\ref{perfect-lemma-affine-compare-bounded}.
We will show that $\omega_A^\bullet$ is a dualizing complex for $A$
which will finish the proof.

\medskip\noindent
Since $X = \bigcup U_i$ is an open covering, we can find a standard
open covering $U = D(f_1) \cup \ldots \cup D(f_m)$ such that
each $D(f_j)$ is a standard open in one of the affine opens $U_i$, see
Schemes, Lemma \ref{schemes-lemma-standard-open-two-affines}.
Say $D(f_j) = D(g_j)$ for $g_j \in A_{i_j}$.
Then $A_{f_j} \cong (A_{i_j})_{g_j}$ and we have
$$
(\omega_A^\bullet)_{f_j} \cong (\omega_i^\bullet)_{g_j}
$$
in the derived category by
Derived Categories of Schemes, Lemma
\ref{perfect-lemma-affine-compare-bounded}.
By Dualizing Complexes, Lemma \ref{dualizing-lemma-dualizing-localize}
we find that
the complex $(\omega_A^\bullet)_{f_j}$ is a dualizing complex over
$A_{f_j}$ for $j = 1, \ldots, m$. This implies that $\omega_A^\bullet$
is dualizing by
Dualizing Complexes, Lemma \ref{dualizing-lemma-dualizing-glue}.
\end{proof}

\end{document}
